% !TeX root = ../../Book5.tex
% 纲要标题：### 27.5 Schubert 演算与四条直线的完整计数
% 纲要位置：几何学计划纲要.md，第 6153 行。
% 正文状态：未撰写；本文件目前是框架。

% 纲要内容开始
% 
% 完整枚举模块固定复数域；一般 Schubert 理论只展开本模型需要的生成元与乘法，不扩展为全部旗簇的表示论。
% 
% * 完全旗、秩条件及 Schubert 子簇的定义
% * 胞腔分解、Chow 群基与特殊 Schubert 类
% * 通用子丛、商丛及低阶 Pieri 乘法
% * $\operatorname{Gr}(2,4)\subset\mathbb P^5$ 的 Plücker 二次方程
% * 与一条固定直线相交的超平面条件
% * 四个一般位置条件的方程消去与交数核验
% 
%   完成 $\int_{\operatorname{Gr}(2,4)}\sigma_1^4=2$，并与 Plücker 二次超曲面和四个超平面条件的计算互相核验：三维复射影空间中与四条一般位置直线均相交的直线有两条。
% 
% * 一般位置、横截性与退化时的重数
% 
%   给出一个一般位置实例验证两解不同，并说明退化时计重的意义；实数域下的实解数不能直接由复交数替代。
% 
% * 与 II.21.5 的胞腔、杯积及复定向比较
% 
% 纲要内容结束

\section{Schubert 演算与四条直线的完整计数}
\label{g5:sec:2705}

待撰写

\subsection{完全旗、秩条件及 Schubert 子簇的定义}
\label{g5:subsec:2705:01}
待撰写

\subsection{胞腔分解、Chow 群基与特殊 Schubert 类}
\label{g5:subsec:2705:02}
待撰写

\subsection{通用子丛、商丛及低阶 Pieri 乘法}
\label{g5:subsec:2705:03}
待撰写

\subsection{\texorpdfstring{\(\operatorname{Gr}(2,4)\subset\mathbb P^5\) 的 Plücker 二次方程}{Gr(2,4)⊂P⁵ 的 Plücker 二次方程}}
\label{g5:subsec:2705:04}
待撰写

\subsection{与一条固定直线相交的超平面条件}
\label{g5:subsec:2705:05}
待撰写

\subsection{四个一般位置条件的方程消去与交数核验}
\label{g5:subsec:2705:06}
待撰写

\subsection{一般位置、横截性与退化时的重数}
\label{g5:subsec:2705:07}
待撰写

\subsection{与 II.21.5 的胞腔、杯积及复定向比较}
\label{g5:subsec:2705:08}
待撰写
% 纲要条目排为编号小节标题，并进入目录与书签；说明段仍保留于纲要注释中。
